\clearpage
\VerdierPhysicalPageAnchor{0233}%
\begingroup\small\setlength{\parskip}{0pt}\setlength{\emergencystretch}{1.5em}
\setlength{\abovedisplayskip}{4pt plus 2pt minus 2pt}%
\setlength{\belowdisplayskip}{4pt plus 2pt minus 2pt}%

\VerdierSemanticAnchor{ch4d-4-1-8}
\noindent\textbf{4.1.8.}\quad
$X$ が単純複体であるとき、第一次数でフィルター付けされた複体 $X$ によって
定まる、$K(\mathcal A)$ に値をとるスペクトル対象 (4.1.7) は、複体 $X$
に付随する、$K(\mathcal A)$ に値をとる第一スペクトル対象にほかならない。

\VerdierSemanticAnchor{ch4d-4-1-8-degree-support}
同様に、$X$ を $n$-重複複体、$m$ を整数、$1\leq m\leq n$ とする。
$X$ の成分 $X^\sigma$ が、次が成り立つとき零であると仮定する：
\[
\sigma-\sigma_m e_m\neq0.
\]
ここで $\sigma_m$ は第 $m$ 座標を表す。このとき、第 $m$ 次数でフィルター
付けされた複体 $X$ によって定まるスペクトル対象は、単純複体 $\int_m X$
の第一標準スペクトル対象に等しいことが確かめられる。

\VerdierSemanticAnchor{ch4d-4-1-8-derived-image}
最後に、圏 $\mathcal A$ がアーベル圏であるとき、標準関手
\[
K(\mathcal A)\longrightarrow D(\mathcal A)
\]
は、第 $m$ 次数でフィルター付けされた $n$-重複複体 $X$ によって定まる、
$K(\mathcal A)$ に値をとるスペクトル対象を、$D(\mathcal A)$ に値をとる
スペクトル対象へ送る。こうして得られるスペクトル対象を、\emph{第 $m$
次数でフィルター付けされた $X$ によって定まる、$D(\mathcal A)$ に値をとる
スペクトル対象}という。混同のおそれがないとき、「$D(\mathcal A)$ に値を
とる」という句を省略する。

\VerdierSemanticAnchor{ch4d-4-1-9}
\medskip\noindent\textbf{4.1.9.}\quad
$X^{\bullet\bullet}$ を $\mathcal A$ の二重複体とする。

\VerdierSemanticAnchor{ch4d-4-1-9-row-complexes}
各整数 $q$ に対し、$X^{q,\bullet}$（それぞれ $X^{\bullet,q}$）を単純複体
\[
\begin{aligned}
X^{q,\bullet}
&=\cdots\xrightarrow{d_2^{q,i-1}}X^{q,i}
  \xrightarrow{d_2^{q,i}}X^{q,i+1}
  \xrightarrow{d_2^{q,i+1}}\cdots,\\[3pt]
(\text{resp.}\quad X^{\bullet,q}
&=\cdots\longrightarrow X^{i,q}
  \xrightarrow{(-1)^q d_1^{i,q}}X^{i+1,q}
  \xrightarrow{(-1)^q d_1^{i+1,q}}X^{i+2,q}
  \longrightarrow\cdots\quad),
\end{aligned}
\]
とする。その次数 $i$ の成分は $X^{q,i}$（それぞれ $X^{i,q}$）である。

\VerdierSemanticAnchor{ch4d-4-1-9-direction-morphisms}
$X^{\bullet\bullet}$ の第一方向（それぞれ第二方向）の微分は、射
（$K(\mathcal A)$ における）
\[
\begin{aligned}
d_1^{q,\bullet}&\colon X^{q,\bullet}\longrightarrow X^{q+1,\bullet}\\
(\text{resp.}\quad
d_2^{\bullet,q}&\colon X^{\bullet,q}\longrightarrow X^{\bullet,q+1}\quad),
\end{aligned}
\]
を誘導する。

\VerdierSemanticAnchor{ch4d-4-1-9-components}
その成分は
\[
\begin{aligned}
(d_1^{q,\bullet})^i&=d_1^{q,i}\\
(\text{resp.}\quad(d_2^{\bullet,q})^i&=(-1)^i d_2^{i,q}\quad).
\end{aligned}
\]
である。
\endgroup
